\documentclass[10pt,a4paper]{amsart}
\usepackage[margin=19mm]{geometry}
\usepackage{amsmath,amssymb,mathtools}
\usepackage[scr=rsfs,cal=euler]{mathalfa}
\usepackage{xfrac}
\usepackage[unicode,hidelinks,bookmarks=false]{hyperref}
\newcommand{\ie}{i.e.\ }
\newcommand{\cf}{cf.\ }
\newcommand{\resp}{respectively\ }
\newcommand{\Subsection}{\subsection}
\providecommand{\nobreakdash}{\nobreak\hbox{-}}
\setlength{\emergencystretch}{2em}
\multlinegap0pt
\usepackage{bm}
\usepackage{bbm}
\usepackage{dsfont}
\usepackage{xspace}
\usepackage{cancel}
\usepackage{mathtools}
\usepackage{amsthm}
\makeatletter
\@ifundefined{theorem}{\newtheorem{theorem}{Theorem}}{}
\@ifundefined{lemma}{\newtheorem{lemma}[theorem]{Lemma}}{}
\makeatother
\newcommand{\Ric}{\operatorname{Ric}}
\newcommand{\dd}{\,\mathrm{d}}
\title[Weighted volume monotonicity and isoperimetric comparison un]{Weighted volume monotonicity and isoperimetric comparison under a lower Ricci curvature bound}
\author{Jiewon Park, Keomkyo Seo}
\date{Source: arXiv:2609.17055v1 (CC BY 4.0)}
\begin{document}
\maketitle

\noindent\emph{Source and licence.} Weighted volume monotonicity and isoperimetric comparison under a lower Ricci curvature bound. Version arXiv:2609.17055v1 (CC BY 4.0), distributed under the Creative Commons Attribution 4.0 International licence, as recorded in the arXiv licence metadata for this exact version and shown on the abstract page https://arxiv.org/abs/2609.17055v1. Full rights notice: licenses/arxiv-2609.17055-CC-BY-4.0.txt.\par\smallskip

\noindent\textit{Editorial scope.} Counted unit: the Theorem whose source label is \texttt{thm:main} in the article (source0.tex lines 194--199, source label \texttt{thm:main}), delivered together with the article's own complete original proof of it (source0.tex lines 201--320, one \texttt{proof} environment of 120 lines). No mathematics is written, repaired or re-derived: statement and proof are the article's own text line for line. Required context retained uncounted: lem:area-comparison (lines 151--153). Every definition, display and label of those lines is retained unchanged. Omitted from this package and not redistributed: 1--150, 154--193, 200--200, 321--1035. Nothing inside a retained excerpt is omitted or altered: every retained body is delivered exactly as the article's lines are written, so no omission receipt of any kind accompanies this package. Citations: the retained excerpts carry no citation command at all, so no bibliography entry of the article is transcribed and no reference is redistributed. Reference adaptation: none was needed and none was made; every reference inside the delivered unit resolves inside it, so no unresolved reference or citation is produced. Printed slips kept literal: no printed word, symbol, hypothesis or number of the counted statement or of its proof was altered. Layout and class: the article's own class is \texttt{article} with packages including amsmath, amssymb, amsthm, mathtools, color, geometry, enumitem; the wrapper is the lane's pinned \texttt{amsart} editorial wrapper at 10pt on A4, which restates the theorem-like environments the retained text needs and adds the article's own macro definitions that the retained text needs (\texttt{\textbackslash Ric}, \texttt{\textbackslash dd}), with extra wrapper packages (bm, bbm, dsfont, xspace, cancel, mathtools). The delivered numbering is the unit's own. Assets: no retained span contains a figure environment, a graphics-inclusion command, a PDF-page inclusion, a table or a TikZ picture; no image file is copied into this package, the package declares no asset dependency and no image file is redistributed with it. Licence: version arXiv:2609.17055v1 of this article is distributed under the Creative Commons Attribution 4.0 International licence, as recorded in the arXiv licence metadata for this exact version and shown on the abstract page https://arxiv.org/abs/2609.17055v1. A full rights notice is delivered at licenses/arxiv-2609.17055-CC-BY-4.0.txt. Characters: every retained excerpt is pure ASCII, so the delivered engine, XeLaTeX with its default Latin Modern text font, reports no missing character.
\begin{lemma} \label{lem:area-comparison}
Let \((M^n,g)\) be an $n(\geq 2)$-dimensional complete Riemannian manifold satisfying $\Ric_M\geq(n-1)kg$ for \(k\in\mathbb{R}\). Then, for $p\in M$, the function $\frac{A(p,r)}{A_k (r)}$ is nonincreasing on \((0,R_k)\), where $A(p,r)=\mathcal H_g^{n-1}(S_p(r))$. In particular, $A(p,r)\leq A_k(r)$ for every $r\in(0,R_k)$. Moreover, for any $r_0\in(0,R_k)$, equality holds if and only if $B_p(r_0)$ is isometric to the geodesic ball of radius $r_0$ in $\mathbb{M}_k^n$.
\end{lemma}

\begin{theorem} \label{thm:main}
Let $(M^n,g)$ be an $n(\geq 2)$-dimensional complete Riemannian manifold satisfying $\Ric_M\geq(n-1)kg$ for \(k\in\mathbb{R}\). For \(p\in M\) and  \(0<R<R_k\), let $f,h\colon(0,R)\rightarrow(0,\infty)$ be positive smooth functions such that 
the functions $f(t)A(p,t), f(t)s_k(t)^{n-1}, h(t)s_k(t)^{n-1}$ are integrable on \((0,R_k)\). Assume that $H(t):=\frac{h(t)}{f(t)}$ is nondecreasing on \((0,R_k)\). Then the function 
$$\Phi_{f,h}(p,r) := \frac{V_f(p,r)}{V_{k,h}(r)}$$
 is nonincreasing in $r \in (0,R_k)$. Moreover, $\Phi_{f,h}(p,r_1)=\Phi_{f,h}(p,r_2)$ holds for some \(0<r_1<r_2<R_k\) if and only if \(H\) is constant on \((0,r_2)\) and $B_p(r_2)$ is isometric to the geodesic ball of radius $r_2$ in $\mathbb M_k^n$.
\end{theorem}

\begin{proof}
By the coarea formula,
\[
V_f(p,r)
=
\int_0^r f(t)A(p,t)\dd t.
\]
Define
\[
q_p(t):=\frac{A(p,t)}{s_k(t)^{n-1}}.
\]
Since $q_p$ is nonincreasing by Lemma~\ref{lem:area-comparison}, we see that
\begin{align}
V_f(p,r)
&=
\int_0^r f(t)s_k(t)^{n-1}q_p(t)\dd t
\nonumber\\
&\geq
q_p(r)
\int_0^r f(t)s_k(t)^{n-1}\dd t.
\label{eq:weighted-lower-bound}
\end{align}
The coarea formula yields at points of differentiability,
\begin{align*}
\partial_r V_f(p,r)&=A_f(p,r)=f(r)A(p,r), \\
V_{k,h}'(r)&=A_{k,h}(r) =\alpha_{n-1}h(r)s_k(r)^{n-1},
\end{align*}
which shows that
\begin{align*}
\partial_r\Phi_{f,h}(p,r)
&=
\frac{
 A_f(p,r)V_{k,h}(r)
 -A_{k,h}(r)V_f(p,r)
}{
 V_{k,h}(r)^2
}.
\end{align*}
Using \eqref{eq:weighted-lower-bound}, we obtain
\begin{align}
\partial_r\Phi_{f,h}(p,r)
&\leq
\frac{\alpha_{n-1}A(p,r)}{V_{k,h}(r)^2}
\left[
 f(r)\int_0^r h(t)s_k(t)^{n-1}\dd t
 -
 h(r)\int_0^r f(t)s_k(t)^{n-1}\dd t
\right]
\nonumber\\
&=
\frac{\alpha_{n-1}A(p,r)}{V_{k,h}(r)^2}
\int_0^r
 f(r)f(t)s_k(t)^{n-1}
 \bigl[H(t)-H(r)\bigr]\dd t
\nonumber\\
&\leq0,
\label{eq:original-proof-chain}
\end{align}
since $H$ is nondecreasing. This proves the desired monotonicity.

Now consider the case of equality. Suppose first that $\Phi_{f,h}(p,r_1)=\Phi_{f,h}(p,r_2)$ for some \(0<r_1<r_2<R_k\). Since \(\Phi_{f,h}(p,\cdot)\) is nonincreasing, it is constant on \([r_1,r_2]\). Hence we see that 
$$\partial_r\Phi_{f,h}(p,r)=0$$
for almost every \(r\in(r_1,r_2)\). It follows from \eqref{eq:original-proof-chain} that
$$H(t)=H(r)$$
for \(t\in(0,r)\). Since this holds for almost every $r\in(r_1,r_2)$ and $H$ is continuous, we conclude that $H$ is constant on $(0,r_2)$.

Since equality has to hold in (1) for every $r \in (r_1, r_2)$ and $t \in (0,r)$, we conclude that
$$
q_p(t)=q_p(r).
$$
Since $q_p$ is nonincreasing, it is constant on \((0,r)\). Together with
\[
\lim_{t\to0^+}q_p(t)=\alpha_{n-1},
\]
this gives
\begin{equation*}
A(p,t)
=
\alpha_{n-1}s_k(t)^{n-1}
=
A_k(t)
\qquad\text{for every }t\in(0,r_2).
\end{equation*}
By Lemma \ref{lem:area-comparison}, we see that $B_p(r_2)$ is isometric to the geodesic ball of radius $r_2$ in $\mathbb{M}_k^n$.









Conversely, suppose that \(H\equiv c\) on \((0,r_2)\) and that \(B_p(r_2)\) is isometric to the geodesic ball of radius $r_2$ in $\mathbb{M}_k^n$. Then, for every \(0<r\leq r_2\),
\[
V_f(p,r)
=
\alpha_{n-1}\int_0^r f(t)s_k(t)^{n-1}\dd t
\]
and
\[
V_{k,h}(r)
=
\alpha_{n-1}\int_0^r c f(t)s_k(t)^{n-1}\dd t
=
cV_f(p,r).
\]
Hence
\[
\Phi_{f,h}(p,r)=\frac1c
\qquad (0<r\leq r_2),
\]
and in particular
\[
\Phi_{f,h}(p,r_1)=\Phi_{f,h}(p,r_2),
\]
which completes the proof.


\end{proof}

\end{document}
